% 2010 年陕西卷理科第 7 题：某空间几何体的三视图。
% 来源 PDF：普通高考/2010/2010陕西理.pdf，SHA-256 8406e303e7aa1ac7813893781c903a7d9c904530028e6fd1c06062d1c32bd645。
% 原页：第 1 页（A3 三栏）第 2 栏上部，书页码 623。
% 重绘依据：照原页——上方为矩形主视图（宽标 √2、高标 1），其右上角右侧另有直角三角形左视图
%   （竖直边在左、斜边自左上落向右下、底边水平），下方为正方形俯视图（右边标 √2）。
%   三个视图按原页的上下、左右相对位置摆放；尺寸线用细线加双向箭头，标注照录。
\begin{tikzpicture}[semithick, x=1cm, y=1cm,
  dim/.style={<->, >=stealth, thin},
  ext/.style={thin}]

  % 主视图：宽 √2≈1.90，高 1≈1.34
  \draw (0,2.60) rectangle (1.90,3.94);
  \node at (0.95,2.35) {主视图};

  % 俯视图：边长 1.90 的正方形
  \draw (0,-0.20) rectangle (1.90,1.70);
  \node at (0.95,-0.45) {俯视图};

  % 主视图宽度尺寸线（√2）
  \draw[ext] (0,3.94) -- (0,4.32);
  \draw[ext] (1.90,3.94) -- (1.90,4.32);
  \draw[dim] (0,4.18) -- (1.90,4.18);
  \node[above] at (0.95,4.18) {$\sqrt2$};

  % 主视图高度尺寸线（1）
  \draw[ext] (1.90,2.60) -- (2.26,2.60);
  \draw[ext] (1.90,3.94) -- (2.26,3.94);
  \draw[dim] (2.12,2.60) -- (2.12,3.94);
  \node[right] at (2.28,3.27) {$1$};

  % 俯视图高度尺寸线（√2）
  \draw[ext] (1.90,-0.20) -- (2.26,-0.20);
  \draw[ext] (1.90,1.70) -- (2.26,1.70);
  \draw[dim] (2.12,-0.20) -- (2.12,1.70);
  \node[right] at (2.28,0.75) {$\sqrt2$};

  % 左视图：直角三角形，竖直直角边在左，斜边自左上落向右下
  \draw (3.60,2.60) -- (3.60,3.94) -- (5.50,2.60) -- cycle;
  \node at (4.55,2.35) {左视图};
\end{tikzpicture}
